% ECONOMICS_PROBLEM: {"id": "D44-238", "title": "Optimal mechanisms for one unit-demand buyer", "jel_code": "D44", "source_id": "OP-0173"}
% !TeX program = xelatex
\documentclass[11pt,a4paper]{article}
\usepackage[margin=23mm]{geometry}
\usepackage{fontspec}
\setmainfont{Latin Modern Roman}
\usepackage{amsmath,amssymb,mathtools}
\usepackage{xcolor,enumitem,booktabs,longtable,array,xurl,hyperref}
\definecolor{muted}{HTML}{53636C}
\hypersetup{colorlinks=true,urlcolor=blue,linkcolor=blue}
\setlength{\parindent}{0pt}
\setlength{\parskip}{5pt}
\newcommand{\R}{\mathbb R}
\newcommand{\E}{\mathbb E}
\newcommand{\N}{\mathbb N}
\newcommand{\ind}{\mathbf 1}
\newcommand{\Var}{\operatorname{Var}}
\newcommand{\Cov}{\operatorname{Cov}}
\newcommand{\supp}{\operatorname{supp}}
\DeclareMathOperator*{\argmin}{arg\,min}
\DeclareMathOperator*{\argmax}{arg\,max}
\newcommand{\entrymeta}[3]{{\sffamily\small\color{muted}\textbf{\texttt{#1}}\quad|\quad #2\par #3\par}}
\newcommand{\Problemref}[1]{\texttt{#1}}
\newcommand{\JELref}[2]{\texttt{#2} (JEL \texttt{#1})}
\begin{document}
\section*{Optimal mechanisms for one unit-demand buyer}
\textbf{Problem ID:} \texttt{D44-238}\quad \textbf{JEL:} \texttt{D44}\par
% BEGIN ECONOMICS STATEMENT
\label{problem:OP-0173}
% GAME_THEORY_PROBLEM: {"local_id": "GT-043", "original_number": 43, "kind": "R", "entry_kind": "statement_only", "published": false, "review_required": true, "category": "game_theory"}
\label{game:GT-043}
{\sffamily\small\color{muted}
\textbf{GT-043} | Mechanism design and auctions | \textbf{R}: published optimal-mechanism agenda, with a precise model.
\par}
\subsection*{Definitions and mathematical statement}
Let $m\geq2$ and $V=\prod_{j=1}^m[a_j,b_j]\subset\mathbb R_{\geq0}^m$. The buyer's values $v_j$ are independent with specified densities $f_j$, and the buyer consumes at most one item. A mechanism consists of measurable $x:V\to\mathbb R_{\geq0}^m$ with $\sum_jx_j(v)\leq1$ and $p:V\to\mathbb R_{\geq0}$. Its utility is $v\cdot x(z)-p(z)$ when type $v$ reports $z$. Impose, for every $v,z\in V$,
\[
 v\cdot x(v)-p(v)\geq v\cdot x(z)-p(z),\qquad v\cdot x(v)-p(v)\geq0.
\]
Determine
\[
 R(F)=\sup_{(x,p)\text{ satisfying these conditions}}\int_Vp(v)\prod_jf_j(v_j)\,dv,
\]
and give optimality criteria and constructions for general $F$. Equivalently, a menu is a set of pairs $(q,t)$ with $q\geq0$, $\sum_jq_j\leq1$, including $(0,0)$; a buyer selects an entry maximizing $v\cdot q-t$.
\subsection*{Short English statement}
Characterize the revenue-optimal menu of lotteries for a buyer who wants only one item.
\subsection*{Variants and scope boundaries}
Restricting to deterministic item prices excludes lotteries and can lower revenue. The source solves important distribution families; those theorems must not be listed as open. An ``exact algorithm'' on arbitrary real densities has no defined bit-complexity meaning without a representation and output specification.
\subsection*{Source relationship and status}
The paper's general agenda motivates this formulation; it does not assert that arbitrary densities admit finite closed-form menus. Its Section 2 supplies the lottery/IC/IR model, and its particular solved families delimit the remaining general characterization.
\subsection*{References}
\begingroup\footnotesize\raggedright
\textbf{[S1]} Kento Hashimoto, Keita Kuwahara, and Reo Nonaka. \emph{Selling Multiple Items to a Unit-Demand Buyer via Automated Mechanism Design}. arXiv:2502.10086v2, paper retrieved with a 2026 manuscript date. \textit{Verified locator:} introduction; Section 2, Model; Theorems 1--2 for solved special families. \url{https://arxiv.org/html/2502.10086v2}
\par\endgroup

\subsection*{Common conventions: Source mathematical conventions}

\textbf{Finite games.} For a finite set $A$, $\Delta(A)$ is its probability
simplex. Players choose independent mixed strategies in Nash equilibrium;
correlation is allowed only when explicitly part of the solution concept.
In a game with payoff functions $u_i$ and strategy profile $x$, an additive
$\varepsilon$ Nash equilibrium satisfies
\[
 u_i(x)\geq u_i(x_i',x_{-i})-\varepsilon
 \quad\text{for every player $i$ and every admissible deviation $x_i'$}.
\]
For numerical approximation, payoffs are normalized as locally specified.
An $\varepsilon$ payoff residual is distinct from distance at most
$\varepsilon$ to the set of exact equilibria. Point convergence, convergence
of empirical play, and convergence in distance to an equilibrium set are
also distinct.

\textbf{Computational representations.} Unless stated otherwise, a complexity
question takes rational data with binary encoding; polynomial time is measured
in total bit length. A fixed-accuracy polynomial algorithm is not necessarily
polynomial in $1/\varepsilon$, and neither is necessarily polynomial in
$\log(1/\varepsilon)$. Succinct games and value, demand, or sampling oracles
must declare their interfaces. Exact real solutions require an explicit
representation; an unspecified list of arbitrary real numbers is not a
finite computational output. A claimed algorithmic barrier is meaningful
only for its specified model and reductions.

\textbf{Dynamic games.} Histories, information, observation, and allowable
randomization are part of the model. A stationary strategy depends only on
the current state; a positional strategy in a deterministic graph game
chooses one action at each controlled vertex. Nash equilibrium at an initial
state and equilibrium after every history are different requirements.
An expectation of a pathwise $\liminf$ payoff, a $\liminf$ of expected averages,
a limiting expected average, and a uniform finite-horizon equilibrium are
different criteria. None is silently substituted for another.

\textbf{Allocation and mechanisms.} A complete allocation assigns every item
exactly once unless otherwise stated. Zero-valued goods, negative values,
capacity constraints, feasibility, lotteries, and copies of goods affect
fairness definitions and are specified locally. Dominant-strategy incentive
compatibility, Bayesian incentive compatibility, universal truthfulness,
and truthfulness in expectation impose different inequalities. Whenever
an objective ratio has zero denominator, its convention is stated or those
instances are excluded.

\textbf{Infinite-dimensional models.} Expectations, integrals, stochastic
equations, and optimization objectives require their local regularity,
integrability, filtrations, and admissible domains. A backward--forward
mean-field system is a boundary-value problem; it is not an initial-value
semiflow without an additional construction. Long-time, large-population,
and vanishing-noise limits require an order of limits and a convergence
topology. If a minimum need not be attained, the objective is its infimum
or supremum and the approximation requirement is stated separately.
% END ECONOMICS STATEMENT
\end{document}
